% answers_short.tex —— 2013 年·江苏卷 · 简短答案速查
% 由 tools/make_short_answers.py 自 answers.tex 的 \gksolhead/\gkanswer 提取，
% 并与 2013/2013_数学真题_江苏_答案要点.md 逐题比对后生成，勿手改（改源 MD 后重跑）。

\begin{gknotice}
\textbf{使用说明}：本节只给最终答案便于快速核对，完整推导见解析版。答案为编写者独立推导，\textbf{非官方标准答案}。
\end{gknotice}

\gkpart{填空题}
\begin{itemize}[leftmargin=2.6em]
  \item[\textbf{1.}] $\pi$。最小正周期 $T=\dfrac{2\pi}{2}=\pi$。
  \item[\textbf{2.}] $5$。$z=(2-\mathrm i)^2=3-4\mathrm i$，故 $|z|=\sqrt{3^2+(-4)^2}=5$。
  \item[\textbf{3.}] $y=\pm\dfrac34x$。双曲线 $\dfrac{x^2}{a^2}-\dfrac{y^2}{b^2}=1$ 的渐近线为 $y=\pm\dfrac ba x$，这里 $a=4$，$b=3$。
  \item[\textbf{4.}] $8$。含 $3$ 个元素的集合共有 $2^3=8$ 个子集。
  \item[\textbf{5.}] $3$。循环中 $(n,a)$ 依次为 $(1,2)\to(2,8)\to(3,26)$，此时 $26<20$ 不成立，输出 $n=3$。
  \item[\textbf{6.}] 乙，$2$。甲、乙的平均数均为 $90$，方差分别为 $4$ 和 $2$，故乙较稳定，方差为 $2$。
  \item[\textbf{7.}] $\dfrac{20}{63}$。$m$ 取奇数的概率为 $\dfrac47$，$n$ 取奇数的概率为 $\dfrac59$，故所求为 $\dfrac47\times\dfrac59=\dfrac{20}{63}$。
  \item[\textbf{8.}] $1:24$。$S_{\triangle ADE}=\dfrac14S_{\triangle ABC}$，$F$ 到底面的距离为棱柱高的一半，故 $V_1=\dfrac13\cdot\dfrac S4\cdot\dfrac h2=\dfrac{Sh}{24}=\dfrac{V_2}{24}$。
  \item[\textbf{9.}] $\left[-2,\dfrac12\right]$。切线为 $y=2x-1$，区域顶点为 $(0,0)$，$\left(\dfrac12,0\right)$，$(0,-1)$，$x+2y$ 在这三点分别取 $0$，$\dfrac12$，$-2$。
  \item[\textbf{10.}] $\dfrac12$。$\overrightarrow{DE}=\dfrac12\overrightarrow{AB}+\dfrac23\overrightarrow{BC}=-\dfrac16\overrightarrow{AB}+\dfrac23\overrightarrow{AC}$，故 $\lambda_1+\lambda_2=\dfrac12$。
  \item[\textbf{11.}] $(-5,0)\cup(5,+\infty)$。$x>0$ 时 $x^2-4x>x$ 得 $x>5$；$x<0$ 时 $f(x)=-x^2-4x>x$ 得 $-5<x<0$；$x=0$ 不满足。
  \item[\textbf{12.}] $\dfrac{\sqrt3}{3}$。由 $d_1=\dfrac{bc}{\sqrt{b^2+c^2}}$，$d_2=\dfrac{b^2}{c}$ 及 $d_2=\sqrt6d_1$ 得 $b^2(b^2+c^2)=6c^4$，代入 $b^2=a^2-c^2$ 得 $6e^4+e^2-1=0$，故 $e^2=\dfrac13$。
  \item[\textbf{13.}] $-1$ 或 $\sqrt{10}$。令 $u=x+\dfrac1x\geqslant2$，则 $AP^2=(u-a)^2+a^2-2$；$a\geqslant2$ 时最小值 $a^2-2=8$ 得 $a=\sqrt{10}$，$a<2$ 时最小值 $2(a-1)^2=8$ 得 $a=-1$。
  \item[\textbf{14.}] $12$。公比 $q=2$，$a_n=2^{n-6}$，所给不等式等价于 $\dfrac{2^n-1}{32}>2^{\frac{n(n-11)}2}$；$n=12$ 成立，$n\geqslant13$ 时右边已大于左边。
\end{itemize}

\gkpart{解答题}
\begin{itemize}[leftmargin=2.6em]
  \item[\textbf{15.}] (1) 证明：$|\boldsymbol a-\boldsymbol b|^2=2-2\cos(\alpha-\beta)=2$，得 $\cos(\alpha-\beta)=0$，故 $\boldsymbol a\cdot\boldsymbol b=0$，即 $\boldsymbol a\perp\boldsymbol b$；(2) $\alpha=\dfrac{5\pi}{6}$，$\beta=\dfrac{\pi}{6}$
  \item[\textbf{16.}] (1) 证明：$AS=AB$ 且 $AF\perp SB$，故 $F$ 为 $SB$ 中点，于是 $EF\parallel AB$，$EG\parallel AC$，所以平面 $EFG\parallel$ 平面 $ABC$；(2) 证明：由平面 $SAB\perp$ 平面 $SBC$ 及 $AF\perp SB$ 得 $AF\perp$ 平面 $SBC$，故 $AF\perp BC$；又 $AB\perp BC$，所以 $BC\perp$ 平面 $SAB$，从而 $BC\perp SA$
  \item[\textbf{17.}] (1) $y=3$ 与 $3x+4y-12=0$。圆心为 $(3,2)$，设切线 $y=kx+3$，由圆心到切线距离为 $1$ 得 $k=0$ 或 $-\dfrac34$；(2) $\left[0,\dfrac{12}{5}\right]$。$MA=2MO$ 化为 $x^2+(y+1)^2=4$，两圆有公共点需圆心距不超过 $3$，即 $5a^2-12a\leqslant0$
  \item[\textbf{18.}] (1) $1040\ \mathrm m$。$\sin A=\dfrac5{13}$，$\sin C=\dfrac45$，$\sin B=\dfrac{63}{65}$，由正弦定理得 $AB=1260\cdot\dfrac{\sin C}{\sin B}=1040$；(2) $\dfrac{35}{37}\ \mathrm{min}$。设乙出发 $t$ 分钟后，距离平方为 $7400t^2-14000t+10000$，在 $t=\dfrac{35}{37}$ 处最小；(3) $\left[\dfrac{1250}{43},\dfrac{625}{14}\right]\ \mathrm{m/min}$。甲到 $C$ 用 $25.2\ \mathrm{min}$，乙到 $C$ 用 $11+\dfrac{500}{v}\ \mathrm{min}$，由等待时间不超过 $3$ 分钟得 $\left|14.2-\dfrac{500}{v}\right|\leqslant3$
  \item[\textbf{19.}] (1) 证明：$c=0$ 时 $b_n=a+\dfrac{(n-1)d}{2}$，由 $b_2^2=b_1b_4$ 得 $d=2a$，故 $S_n=n^2a$，从而 $S_{nk}=n^2S_k$；(2) 证明：设 $b_n=pn+q$，与 $b_n=\dfrac{nS_n}{n^2+c}$ 比较系数，由 $n^3$ 项得 $p=\dfrac d2\neq0$，由 $n$ 项得 $pc=0$，故 $c=0$
  \item[\textbf{20.}] (1) $(e,+\infty)$。$f'(x)=\dfrac1x-a\leqslant0$ 在 $(1,+\infty)$ 上恒成立需 $a\geqslant1$；$g'(x)=\mathrm e^x-a$ 在 $(1,+\infty)$ 上有零点且变号需 $a>e$；(2) 当 $a\leqslant0$ 时 $1$ 个零点；当 $0<a<\dfrac1e$ 时 $2$ 个零点；当 $a=\dfrac1e$ 时 $1$ 个零点
\end{itemize}

\gkpart{附加题}
\begin{itemize}[leftmargin=2.6em]
  \item[\textbf{21.}] (1) 证明：设 $OC=r$，由坐标计算得 $AC=\dfrac{8r}{3}$，切线长 $AD=\dfrac{4r}{3}$，故 $AC=2AD$；(2) $\begin{bmatrix}-1&-2\\0&3\end{bmatrix}$；(3) $l:2x-y-2=0$，$C:y^2=2x$，公共点为 $(2,2)$，$\left(\dfrac12,-1\right)$；(4) 证明：$2a^3-b^3-(2ab^2-a^2b)=(2a+b)(a+b)(a-b)\geqslant0$
  \item[\textbf{22.}] (1) $\dfrac{3\sqrt{10}}{10}$。以 $A$ 为原点建立空间直角坐标系，$\overrightarrow{A_1B}=(2,0,-4)$，$\overrightarrow{C_1D}=(1,-1,-4)$，夹角余弦为 $\dfrac{3}{\sqrt{10}}$；(2) $\dfrac{\sqrt5}{3}$。平面 $ADC_1$ 的法向量取 $(2,-2,1)$，平面 $ABA_1$ 的法向量为 $(0,1,0)$，两法向量夹角余弦为 $\dfrac23$，故二面角正弦为 $\dfrac{\sqrt5}{3}$
  \item[\textbf{23.}] (1) $5$。$P_{11}=\{1,4,5,6,11\}$；(2) $1008$。奇数编号的块中所有项都满足，前 $61$ 块贡献 $31^2=961$ 项，第 $63$ 块中 $n=1954,\cdots,2000$ 贡献 $47$ 项
\end{itemize}
